\clearpage
\section{Differential connections and colored
forms}\label{g-differential-connections-and-colored-forms}

Technical derivation for review. Equation and subsection numbers in this
appendix are local to the source derivation. Historical local labels are
retained, including numbering gaps or nonsequential labels. The
accompanying source notes retain the original expressions for
comparison.

\subsection{1. Setting and result}\label{g-setting-and-result}

Fix the rank and block partition. Assume only

% DISPLAY g.1
% SOURCE alpha_a+alpha_b <= kappa r_ab,     |Y| <= sigma r_*,
% SOURCE r_* = min_e r_e >= 1,             W_2 = sum_e r_e^-2.
\begin{gather*}
\begin{aligned}
& \alpha _{a}+\alpha _{b} \le  \kappa  r_{ab}, \quad  \vert Y\vert  \le  \sigma  r_{*}, \\
& r_{*} = \min _{e} r_{e} \ge  1, \quad  W_{2} = \sum _{e} r_{e}^{-2}.
\end{aligned}\notag
\end{gather*}

Constants depend on fixed rank and fixed sufficiently small upper bounds
for \(\kappa\) and \(\sigma\), never on \(r_{\max}/r_{*}\). Work first
on the compactly supported smooth residual-equivariant core in a
normal-frame chart. No individual block-singlet assumption is imposed.
Let

% DISPLAY g.2
% SOURCE H_exc = sum_e (H_fast,e-E_0,e) >= 0,
% SOURCE lambda(Z) = sigma^2+r_*^-3,
% SOURCE E[v] = (lambda H_exc)[v]+W_2[v].                       (1.1)
\begin{gather*}
\begin{aligned}
& H_{\mathrm{exc}} = \sum _{e} (H_{\mathrm{fast},e}-E_{0,e}) \ge  0, \\
& \lambda (Z) = \sigma ^{2}+r_{*}^{-3}, \\
& E[v] = (\lambda  H_{\mathrm{exc}})[v]+W_{2}[v].
\end{aligned}\tag{1.1}
\end{gather*}

The weights are center-only and the excitation form acts in fast fibers.
For the INTERNAL part of the completion connection write

% DISPLAY g.3
% SOURCE beta_j = b_j(Y,A,Z) dot p_F,
% SOURCE beta = A_s^* B_F p_F,     A_s = B_s G_s^-1/2.
\begin{gather*}
\begin{aligned}
& \beta _{j} = b_{j}(Y,A,Z) \cdot  p_{F}, \\
& \beta  = A_{s}^{*} B_{F} p_{F}, \quad  A_{s} = B_{s} G_{s}^{-1/2}.
\end{aligned}\notag
\end{gather*}

Here \(A_{s}\) is the normalized Gauss slow coefficient, not the
internal matrix tuple. The internal block of \(G_{s}\) is exactly
identity. Define the ordered shifted form inherited from the exact
kinetic identity:

% DISPLAY g.4
% SOURCE H_I^raw(beta)[v]
% SOURCE   = sum_j ||(p_j+beta_j)v||^2+(V_I+Y_I)[v].             (1.2)
\begin{gather*}
\begin{aligned}
& H_{I}^{\mathrm{raw}}(\beta )[v] \\
& = \sum _{j} \Vert (p_{j}+\beta _{j})v\Vert ^{2}+(V_{I}+Y_{I})[v].
\end{aligned}\tag{1.2}
\end{gather*}

For every \(\epsilon\) in (0,1),

% DISPLAY g.5
% SOURCE H_I^raw(beta)[v]
% SOURCE   >= (1-epsilon)H_I[v]
% SOURCE        -C_(N,epsilon)(lambda H_exc)[v]
% SOURCE        -C_(N,epsilon)W_2[v].                         (1.3)
\begin{gather*}
\begin{aligned}
& H_{I}^{\mathrm{raw}}(\beta )[v] \\
& \ge  (1-\epsilon )H_{I}[v] \\
& -C_{N,\epsilon}(\lambda  H_{\mathrm{exc}})[v] \\
& -C_{N,\epsilon}W_{2}[v].
\end{aligned}\tag{1.3}
\end{gather*}

This retains the ENTIRE interacting internal form, including its
original potential and Yukawa terms. There is no \(\epsilon\) times bare
internal kinetic loss and no \(\epsilon  \alpha _{a}\) mass remainder.
No internal gap or physical-sector lower-rank Hardy estimate is used.

The proof gives the actual curvature expectation bound

% DISPLAY g.6
% SOURCE |R_beta[v]| <= C_N(lambda H_exc)[v]+C_N W_2[v].       (1.4)
\begin{gather*}
\begin{aligned}
& \vert R_{\beta}[v]\vert  \le  C_{N}(\lambda  H_{\mathrm{exc}})[v]+C_{N} W_{2}[v].
\end{aligned}\tag{1.4}
\end{gather*}

The stronger possible coefficient \(C_{N} r_{*}^{-3}\) with
\(\sigma ^{2}\) absent is not needed and is not claimed. Fixing
\(\sigma\) sufficiently small and then taking the exterior radius large
makes the excitation coefficient as small as desired. On fast-complement
inputs \(W_{2}\) is also paid by the gap.

\subsection{2. Exact completion}\label{g-exact-completion}

Remove the finite fermionic Gauss source by the separate estimate chosen
for that source. Its estimate is not part of this algebra. Let

% DISPLAY g.7
% SOURCE w=G_s^1/2 p_s v,   a=B_F p_F v,   A_s=B_s G_s^-1/2.
\begin{gather*}
\begin{aligned}
& w=G_{s}^{1/2} p_{s} v, \quad  a=B_{F} p_{F} v, \quad  A_{s}=B_{s} G_{s}^{-1/2}.
\end{aligned}\notag
\end{gather*}

Exactly, as a first-order form identity,

% DISPLAY g.8
% SOURCE ||w||^2+||a+A_s w||^2
% SOURCE   =||w+A_s^*a||^2+||a||^2+||A_s w||^2-||A_s^*a||^2. (2.1)
\begin{gather*}
\begin{aligned}
& \Vert w\Vert ^{2}+\Vert a+A_{s} w\Vert ^{2} \\
& =\Vert w+A_{s}^{*}a\Vert ^{2}+\Vert a\Vert ^{2}+\Vert A_{s} w\Vert ^{2}-\Vert A_{s}^{*}a\Vert ^{2}.
\end{aligned}\tag{2.1}
\end{gather*}

Coefficients multiply the differentiated section in their original
order. No integration by parts and no coefficient derivative occurs in
this completion.

The internal components of \(w+A_{s}^{*}a\) are \(p_{j}+\beta _{j}\).
Its center components use the exact center metric and NORMAL-BUNDLE
covariant momentum. The center square and \(\Vert A_{s} w\Vert ^{2}\)
may be retained or discarded as positive forms. They are not additional
copies of the fast Gauss energy. The negative term is the connection
norm already bounded in the rootwise lemma.

If the half-density was separated as in the intrinsic geometry note, its
scalar \(V_{d}\) must still be included separately. The symmetrization
below neither replaces nor duplicates that calculation.

\subsection{3. Coefficients and differentiated rootwise
bounds}\label{g-coefficients-and-differentiated-rootwise-bounds}

Use intrinsic notation

% DISPLAY g.9
% SOURCE F=E^*B_(Z+A),   C=P_N B_A,   N_Y=P_N B_Y,
% SOURCE I_Y=P_I B_Y,    M=F+L-J_2,   W=I+KK^*.
\begin{gather*}
\begin{aligned}
& F=E^{*}B_{Z+A}, \quad  C=P_{N} B_{A}, \quad  N_{Y}=P_{N} B_{Y}, \\
& I_{Y}=P_{I} B_{Y}, \quad  M=F+L-J_{2}, \quad  W=I+KK^{*}.
\end{aligned}\notag
\end{gather*}

The internal coefficient is, with the common Gauss sign convention,

% DISPLAY g.10
% SOURCE beta_I=I_Y M^-1 W M^-* (C+N_Y)^*p_F.                   (3.1)
\begin{gather*}
\begin{aligned}
& \beta _{I}=I_{Y} M^{-1} W M^{-*} (C+N_{Y})^{*}p_{F}.
\end{aligned}\tag{3.1}
\end{gather*}

All coefficients are real in real trace-orthonormal frames. Normal
frames depend on centers only; an internal derivative differentiates
neither \(p_{F}\) nor the frame. The connection commutes with internal
Clifford matrices and every internal potential-supercharge
multiplication coefficient.

For a unit internal direction \(h\),

% DISPLAY g.11
% SOURCE partial_h M=partial_h F,
% SOURCE partial_h F and partial_h C are pair-block diagonal,
% SOURCE ||partial_h F_e||+||partial_h C_e|| <= C_N,
% SOURCE partial_h W=partial_h I_Y=partial_h N_Y=0.             (3.2)
\begin{gather*}
\begin{aligned}
& \partial _{h} M=\partial _{h} F, \\
& \partial _{h} F \text{ and } \partial _{h} C \text{ are pair-block diagonal }, \\
& \Vert \partial _{h} F_{e}\Vert +\Vert \partial _{h} C_{e}\Vert  \le  C_{N}, \\
& \partial _{h} W=\partial _{h} I_{Y}=\partial _{h} N_{Y}=0.
\end{aligned}\tag{3.2}
\end{gather*}

Only pairs incident to the differentiated block contribute. The coarse
\(r_{*}\) estimate below does not convert any kinetic or Yukawa term, so
it creates no unrelated \(\alpha _{a}/r_{*}\) error.

The cited rootwise connection proof gives

% DISPLAY g.12
% SOURCE sum_j ||beta_jv||^2 <= C_N E[v].                     (3.3)
\begin{gather*}
\begin{aligned}
& \sum _{j} \Vert \beta _{j} v\Vert ^{2} \le  C_{N} E[v].
\end{aligned}\tag{3.3}
\end{gather*}

Its explicit first-order decomposition also gives

% DISPLAY g.13
% SOURCE sum_(j,k)|| (partial_j beta_k)v ||^2
% SOURCE    <= C_N [ (r_*^-2 lambda H_exc)[v]
% SOURCE                  +(r_*^-2 W_2)[v] ].                 (3.4)
\begin{gather*}
\begin{aligned}
& \sum _{j,k}\Vert  (\partial _{j} \beta _{k})v \Vert ^{2} \\
& \le  C_{N} [ (r_{*}^{-2} \lambda  H_{\mathrm{exc}})[v] \\
& +(r_{*}^{-2} W_{2})[v] ].
\end{aligned}\tag{3.4}
\end{gather*}

Here \(\partial _{j} \beta _{k}\) differentiates its coefficients only.
The following proves (3.4); it is not inferred from the abstract bound
(3.3).

\subsubsection{3.1 Proof of the differentiated
estimate}\label{g-proof-of-the-differentiated-estimate}

In the rootwise lemma put

% DISPLAY g.14
% SOURCE d_0=F^-1 C^*p_F,    T=(I+F^-1 P^*)^-1,    P=L-J_2.
\begin{gather*}
\begin{aligned}
& d_{0}=F^{-1} C^{*}p_{F}, \quad  T=(I+F^{-1} P^{*})^{-1}, \quad  P=L-J_{2}.
\end{aligned}\notag
\end{gather*}

The exact decompositions (5)-(7) and (12)-(15) in that lemma express
\(\beta\) as finite sums of bounded multiplication matrices times the
source classes

% DISPLAY g.15
% SOURCE Y_e p_e/r_e,
% SOURCE Y_f p_e/r_g,       when (e,f,g) is a triangle,
% SOURCE Y_f p_e/r_e,       with the input-root denominator,
\begin{gather*}
\begin{aligned}
& Y_{e} p_{e}/r_{e}, \\
& Y_{f} p_{e}/r_{g}, \quad  \text{ when } (e,f,g) \text{ is a triangle }, \\
& Y_{f} p_{e}/r_{e}, \quad  \text{ with the input-root denominator },
\end{aligned}\notag
\end{gather*}

with bounded root-preserving factors between the displayed operators.
Non-root-preserving inverse matrices remain bounded LEFT factors. Thus
no unweighted spectator momentum is hidden.

Differentiate those exact decompositions.
\(Y, r_{e}, W, K, I_{Y}, C_{Y}, L, J_{2}\) and \(g_{C}\) have zero
internal derivatives. The differentiated factors satisfy

% DISPLAY g.16
% SOURCE ||partial_j F_e^-1|| <= C_N/r_e^2,
% SOURCE ||partial_j(F_e^-1 C_e^*)|| <= C_N/r_e,
% SOURCE ||partial_j T||+||partial_j T^*|| <= C_N/r_*.
\begin{gather*}
\begin{aligned}
& \Vert \partial _{j} F_{e}^{-1}\Vert  \le  C_{N}/r_{e}^{2}, \\
& \Vert \partial _{j}(F_{e}^{-1} C_{e}^{*})\Vert  \le  C_{N}/r_{e}, \\
& \Vert \partial _{j} T\Vert +\Vert \partial _{j} T^{*}\Vert  \le  C_{N}/r_{*}.
\end{aligned}\notag
\end{gather*}

The inverse derivative formula and (3.2) prove these inequalities. All
other bounded left factors in the displayed decompositions consequently
have derivative at most \(C_{N}/r_{*}\). Differentiating
\(F_{e}^{-1} C_{e}^{*}\) replaces its \(O(\kappa )\) factor by
\(O(1/r_{e})\); use the source bound with its optional \(\kappa\) gain
omitted.

Every resulting term is therefore the SAME source class with one extra
factor at most \(C_{N}/r_{*}\). Applying the proved same-root and
mixed-root inequalities (2)-(4) of the cited lemma, and finite-sum
Cauchy, proves (3.4). The argument is all-vector, does not bound
internal momenta, and does not require a fast projection to preserve
tube support.

\subsubsection{3.2 Coarse pointwise derivatives for
divergence}\label{g-coarse-pointwise-derivatives-for-divergence}

Only for the scalar divergence correction we need

% DISPLAY g.17
% SOURCE |b| <= C_N sigma(kappa+sigma),
% SOURCE |partial_Y b| <= C_N/r_*,
% SOURCE |partial_A partial_Y b|+|partial_Y^2 b| <= C_N/r_*^2. (3.5)
\begin{gather*}
\begin{aligned}
& \vert b\vert  \le  C_{N} \sigma (\kappa +\sigma ), \\
& \vert \partial _{Y} b\vert  \le  C_{N}/r_{*}, \\
& \vert \partial _{A} \partial _{Y} b\vert +\vert \partial _{Y}^{2} b\vert  \le  C_{N}/r_{*}^{2}.
\end{aligned}\tag{3.5}
\end{gather*}

Indeed the internal normalized slow column is \(O(\sigma )\), the full
fast Gauss column is \(O(\kappa +\sigma )\), and
\(\Vert M^{-1}\Vert \le C_{N}/r_{*}\). First derivatives of \(M, I_{Y}\)
and \(C+N_{Y}\) are \(O(1)\); second derivatives of \(M\) are
\(O(1/r_{*})\), coming from \(J_{2}\). First and second derivatives of
\(W\) and its positive square root are \(O(1/r_{*})\) and
\(O(1/r_{*}^{2})\). Differentiate the product \(B_{I}^{*}B_{F}\) using
the inverse derivative formula. The potentially large \(C\) is always
kept inside \(F^{-1} C^{*}\) or the bounded full fast column, rather
than estimated by \(\kappa  r_{\max}\). This proves (3.5); fixed
coordinate dimensions absorb component sums.

\subsection{4. Exact self-adjoint
symmetrization}\label{g-exact-self-adjoint-symmetrization}

In flat fast-fiber measure, with \(p_{F}=-i \partial _{Y}\), put

% DISPLAY g.18
% SOURCE d_j=div_F b_j,
% SOURCE beta_hat_j=(beta_j+beta_j^*)/2
% SOURCE           =b_j dot p_F-(i/2)d_j,
% SOURCE P_j=p_j+beta_hat_j,
% SOURCE X_j=partial_j+b_j dot partial_Y.
\begin{gather*}
\begin{aligned}
& d_{j}=\operatorname{div} _{F} b_{j}, \\
& \widehat{\beta}_{j}=(\beta _{j}+\beta _{j}^{*})/2 \\
& =b_{j} \cdot  p_{F}-(i/2)d_{j}, \\
& P_{j}=p_{j}+\widehat{\beta}_{j}, \\
& X_{j}=\partial _{j}+b_{j} \cdot  \partial _{Y}.
\end{aligned}\notag
\end{gather*}

Then \(\beta _{j}=\widehat{\beta}_{j}+(i/2)d_{j}\) and
\(P_{j}=-i(X_{j}+d_{j}/2)\) is symmetric on the core. Direct integration
by parts gives the EXACT identity

% DISPLAY g.19
% SOURCE ||(p_j+beta_j)v||^2
% SOURCE   =||P_jv||^2
% SOURCE       + integral [(1/2)X_j d_j+(1/4)d_j^2]|v|^2.     (4.1)
\begin{gather*}
\begin{aligned}
& \Vert (p_{j}+\beta _{j})v\Vert ^{2} \\
& =\Vert P_{j} v\Vert ^{2} \\
& + \int  [(1/2)X_{j} d_{j}+(1/4)d_{j}^{2}]\vert v\vert ^{2}.
\end{aligned}\tag{4.1}
\end{gather*}

The cross term is \((1/2) \int  X_{j}(d_{j})\vert v\vert ^{2}\); div
\(X_{j}=d_{j}\). Thus simply treating \(b \cdot  p_{F}\) as symmetric
would miss a scalar correction.

By (3.5),

% DISPLAY g.20
% SOURCE |d_j| <= C_N/r_*,
% SOURCE |partial_A d_j|+|partial_Y d_j| <= C_N/r_*^2,
% SOURCE sum_j |(1/2)X_j d_j+(1/4)d_j^2| <= C_N W_2.           (4.2)
\begin{gather*}
\begin{aligned}
& \vert d_{j}\vert  \le  C_{N}/r_{*}, \\
& \vert \partial _{A} d_{j}\vert +\vert \partial _{Y} d_{j}\vert  \le  C_{N}/r_{*}^{2}, \\
& \sum _{j} \vert (1/2)X_{j} d_{j}+(1/4)d_{j}^{2}\vert  \le  C_{N} W_{2}.
\end{aligned}\tag{4.2}
\end{gather*}

Hence the full correction is inverse-square multiplication. Also

% DISPLAY g.21
% SOURCE sum_j ||beta_hat_jv||^2 <= C_N E[v],                  (4.3)
\begin{gather*}
\begin{aligned}
& \sum _{j} \Vert \widehat{\beta}_{j} v\Vert ^{2} \le  C_{N} E[v],
\end{aligned}\tag{4.3}
\end{gather*}

and (3.4) holds with \(\widehat{\beta}\) replacing \(\beta\): its extra
multiplication derivative is \(O(r_{*}^{-2})\), whose squared norm is
bounded by \(r_{*}^{-2} W_{2}\).

\subsection{5. Actual curvature
expectation}\label{g-actual-curvature-expectation}

Define the symmetric curvature operator

% DISPLAY g.22
% SOURCE F_jk=i[P_j,P_k]
% SOURCE     =partial_j beta_hat_k-partial_k beta_hat_j
% SOURCE                      +i[beta_hat_j,beta_hat_k].       (5.1)
\begin{gather*}
\begin{aligned}
& F_{jk}=i[P_{j},P_{k}] \\
& =\partial _{j} \widehat{\beta}_{k}-\partial _{k} \widehat{\beta}_{j} \\
& +i[\widehat{\beta}_{j},\widehat{\beta}_{k}].
\end{aligned}\tag{5.1}
\end{gather*}

Equivalently \(F_{jk}\) is the symmetrized momentum of the fast vector
field

% DISPLAY g.23
% SOURCE partial_j b_k-partial_k b_j
% SOURCE   +(b_j dot partial_Y)b_k-(b_k dot partial_Y)b_j.
\begin{gather*}
\begin{aligned}
& \partial _{j} b_{k}-\partial _{k} b_{j} \\
& +(b_{j} \cdot  \partial _{Y})b_{k}-(b_{k} \cdot  \partial _{Y})b_{j}.
\end{aligned}\notag
\end{gather*}

It is first order, not bounded multiplication.

The charge identity contracts (5.1) with fixed bounded internal-Clifford
matrices. If \(L\) is any such matrix, it commutes with
\(\widehat{\beta}\) and its coefficient derivatives. Symmetry gives
exactly

% DISPLAY g.24
% SOURCE |<v,L i[beta_hat_j,beta_hat_k]v>|
% SOURCE    <=2||L|| ||beta_hat_jv|| ||beta_hat_kv||.           (5.2)
\begin{gather*}
\begin{aligned}
& \vert \langle v,L i[\widehat{\beta}_{j},\widehat{\beta}_{k}]v\rangle \vert \\
& \le 2\Vert L\Vert  \Vert \widehat{\beta}_{j} v\Vert  \Vert \widehat{\beta}_{k} v\Vert .
\end{aligned}\tag{5.2}
\end{gather*}

Thus (4.3) pays all beta-beta curvature terms. There is no need to bound
the norm of a fast-differentiated curvature symbol.

For the derivative part, (3.4) and Cauchy give

% DISPLAY g.25
% SOURCE |<v,L(partial_j beta_hat_k)v>|
% SOURCE    <= C_N ||r_*^-1 v|| E[v]^1/2
% SOURCE    <= C_N E[v],                                     (5.3)
\begin{gather*}
\begin{aligned}
& \vert \langle v,L(\partial _{j} \widehat{\beta}_{k})v\rangle \vert \\
& \le  C_{N} \Vert r_{*}^{-1} v\Vert  E[v]^{1/2} \\
& \le  C_{N} E[v],
\end{aligned}\tag{5.3}
\end{gather*}

since \(\Vert r_{*}^{-1}v\Vert ^{2}\le W_{2}[v]\le E[v]\). To justify
the first inequality with varying centers, apply the coefficient
derivative bound at each center and then apply weighted Cauchy over
centers. Summing the finitely many Clifford components proves (1.4).

This is an all-vector curvature EXPECTATION estimate. It derives the
extra coefficient derivative from the actual rootwise expansion, not
from a generic relative-norm assertion. It makes no oscillator-degree
cutoff and loses no internal kinetic energy.

\subsection{6. Averaged sixteen-supercharge
comparison}\label{g-averaged-sixteen-supercharge-comparison}

Normalize the internal charges by

% DISPLAY g.26
% SOURCE H_I[v]=(1/16)sum_alpha ||Q_alpha v||^2,
% SOURCE Q_alpha=sum_j c_(alpha,j)p_j+V_alpha(A).
\begin{gather*}
\begin{aligned}
& H_{I}[v]=(1/16)\sum _{\alpha} \Vert Q_{\alpha} v\Vert ^{2}, \\
& Q_{\alpha}=\sum _{j} c_{\alpha ,j}p_{j}+V_{\alpha}(A).
\end{aligned}\notag
\end{gather*}

The \(c_{\alpha ,j}\) are fixed internal Clifford matrices. The ordinary
colored Gauss terms cancel in the spinor average. An individual
\(Q_{\alpha}^{2}\) must not be identified with \(H_{I}\) on the colored
core.

Replace \(p_{j}\) by \(P_{j}\):

% DISPLAY g.27
% SOURCE Q_alpha^beta=Q_alpha+C_alpha(beta_hat),
% SOURCE C_alpha(beta_hat)=sum_j c_(alpha,j) beta_hat_j.
\begin{gather*}
\begin{aligned}
& Q_{\alpha}^{\beta}=Q_{\alpha}+C_{\alpha}(\widehat{\beta}), \\
& C_{\alpha}(\widehat{\beta})=\sum _{j} c_{\alpha ,j} \widehat{\beta}_{j}.
\end{aligned}\notag
\end{gather*}

The exact averaged expansion is

% DISPLAY g.28
% SOURCE (1/16)sum_alpha ||Q_alpha^beta v||^2
% SOURCE    =sum_j||P_jv||^2+(V_I+Y_I)[v]+R_beta[v].           (6.1)
\begin{gather*}
\begin{aligned}
& (1/16)\sum _{\alpha} \Vert Q_{\alpha}^{\beta} v\Vert ^{2} \\
& =\sum _{j}\Vert P_{j} v\Vert ^{2}+(V_{I}+Y_{I})[v]+R_{\beta}[v].
\end{aligned}\tag{6.1}
\end{gather*}

Here \(R_{\beta}\) is a finite-Clifford linear contraction of
\(F_{jk}\). The symmetric momentum products give the kinetic form,
antisymmetric momentum products give curvature, and

% DISPLAY g.29
% SOURCE [P_j,V_alpha]=[p_j,V_alpha]
\begin{gather*}
\begin{aligned}
& [P_{j},V_{\alpha}]=[p_{j},V_{\alpha}]
\end{aligned}\notag
\end{gather*}

because \(\widehat{\beta}\) is fast-bosonic with scalar bosonic
coefficients. The original potential and Yukawa terms are consequently
unchanged. Spinor-averaged cancellation of the original Gauss term is
unchanged too.

Finite-dimensional Clifford bounds and (4.3) imply

% DISPLAY g.30
% SOURCE (1/16)sum_alpha ||C_alpha(beta_hat)v||^2 <= C_N E[v].  (6.2)
\begin{gather*}
\begin{aligned}
& (1/16)\sum _{\alpha} \Vert C_{\alpha}(\widehat{\beta})v\Vert ^{2} \le  C_{N} E[v].
\end{aligned}\tag{6.2}
\end{gather*}

Young\textquotesingle s inequality on the COMPLETE charge stack now
gives

% DISPLAY g.31
% SOURCE (1/16)sum_alpha ||Q_alpha^beta v||^2
% SOURCE    >=(1-epsilon)H_I[v]-C_N epsilon^-1 E[v].           (6.3)
\begin{gather*}
\begin{aligned}
& (1/16)\sum _{\alpha} \Vert Q_{\alpha}^{\beta} v\Vert ^{2} \\
& \ge (1-\epsilon )H_{I}[v]-C_{N} \epsilon ^{-1} E[v].
\end{aligned}\tag{6.3}
\end{gather*}

Subtract curvature using (1.4), then apply (4.1)-(4.2). This proves
(1.3). The loss is \(\epsilon  H_{I}\), not \(\epsilon  t_{I}\); that is
why no large unrelated internal Yukawa mass is introduced.

\subsection{7. Normal-bundle center connection and chart
changes}\label{g-normal-bundle-center-connection-and-chart-changes}

The normal fast bundle \(N_{Z}\) depends on centers alone. Its
orthogonal connection is already in \(p_{C}\) in the exact intrinsic
formula. In a local frame its center generator may contain unbounded
\(Y \partial _{Y}\) operators. This argument never estimates those by
\(H_{\mathrm{exc}}\).

Internal derivatives act in a center-dependent but A-independent frame,
so (3.2) has no frame derivative. The internal supercharges contain no
center momentum. Neither center normal-bundle curvature nor mixed
center/internal curvature occurs in the INTERNAL square (6.1).

The completed center square remains a genuine positive covariant
derivative form. Any subsequent center Hardy argument must use that
actual connection with its domain justified. It may not replace its
unbounded normal-frame generator by a bounded vacuum Berry coefficient.
That separate comparison is not claimed here.

All identities intertwine under center-dependent orthogonal changes of
fast normal frame and residual gauge changes. \(\beta\) is an actual
vector field on the fast fiber, \(\widehat{\beta}\) is its half-density
action, and curvature is its commutator. The internal comparison is
therefore intrinsic on a fixed projector branch.

\subsection{8. Parameter order and exact
scope}\label{g-parameter-order-and-exact-scope}

Choose the requested internal loss \(\epsilon\) first, then \(\sigma\)
so \(C_{N,\epsilon}\sigma ^{2}\) fits the available fast reserve, then
\(r_{*}\) large enough for \(C_{N,\epsilon}r_{*}^{-3}\) and the
inverse-square multiplier.

On the global fast complement \(H_{\mathrm{exc}}\ge c_{N} r_{*}\) and
\(W_{2}\le m_{N} r_{*}^{-2}\), so

% DISPLAY g.32
% SOURCE W_2[v] <= C_N (r_*^-3 H_exc)[v].
\begin{gather*}
\begin{aligned}
& W_{2}[v] \le  C_{N} (r_{*}^{-3} H_{\mathrm{exc}})[v].
\end{aligned}\notag
\end{gather*}

This final payment is restricted to complementary inputs; (1.3) itself
holds on arbitrary colored inputs. No internal Hamiltonian is inverted.

Still separate are the finite spin source, density multiplier, signed
fast kinetic remainder, potential/Yukawa remainders, and branch/cutoff
compatibility. The present lemma closes the complete internal-form and
curvature step only.
